同事发来一份数据集说明：某列标签有 \(99\%\) 是``正常''；另一列是连续的响应延迟；还有一列缺了大半，需要补一个分布进去。三件事看起来毫无关系，可它们背后是同一个问题——``不确定''这个词能不能变成一个可以加减、可以比较、可以当预算来花的数。

本章给出的答案只有一句话：把一个结果的代价定义成 \(-\log P\)，再取平均。这句话很短，落地时却有四个台阶，而中间那一级最容易踩空。

\begin{center}
\begin{tikzpicture}[x=1cm,y=1cm,font=\small,>=stealth]
  \draw[rounded corners,black!45,fill=blue!5] (0,0) rectangle (3.3,1.4);
  \draw[rounded corners,black!45,fill=blue!5] (3.8,0) rectangle (7.1,1.4);
  \draw[rounded corners,black!45,fill=blue!5] (7.6,0) rectangle (10.9,1.4);
  \draw[rounded corners,black!45,fill=blue!5] (11.4,0) rectangle (14.7,1.4);
  \draw[->,black!60] (3.32,0.7) -- (3.78,0.7);
  \draw[->,black!60] (7.12,0.7) -- (7.58,0.7);
  \draw[->,black!60] (10.92,0.7) -- (11.38,0.7);
  \node[align=center,font=\footnotesize] at (1.65,0.7) {一个结果的代价\\\(-\log P\)};
  \node[align=center,font=\footnotesize] at (5.45,0.7) {取平均得熵\\平均问几个是非题};
  \node[align=center,font=\footnotesize] at (9.25,0.7) {换成连续变量\\换个单位数就变};
  \node[align=center,font=\footnotesize] at (13.05,0.7) {反着用\\由约束定出分布};
  \node[align=center,font=\footnotesize,black!65] at (1.65,-0.55) {自信息};
  \node[align=center,font=\footnotesize,black!65] at (5.45,-0.55) {Shannon 熵};
  \node[align=center,font=\footnotesize,black!65] at (9.25,-0.55) {微分熵};
  \node[align=center,font=\footnotesize,black!65] at (13.05,-0.55) {最大熵原理};
\end{tikzpicture}

\emph{四篇的走向。前两步把不确定性变成数，第三步说明这个数换到连续情形会失去绝对意义，第四步把箭头掉转过来：已知若干约束，反推最不多话的那个分布。}
\end{center}

\subsection{本章想让你改掉的几个判断}

第一，熵不是抽象的``混乱度''，它有一个可以数的读法：\emph{平均需要问多少个是非题，才能把结果定位下来}。这个读法让 \(\log_2\) 和 bit 有了实义，也直接解释了熵为什么就是压缩的下界——一串是非题的答案就是一串比特。语言模型报的困惑度是熵的指数，读作``平均在多少个等可能候选之间犹豫''。

第二，熵只看概率，不看取值。把标签重命名、把数值统统放大一千倍，熵一动不动。所以它适合度量``有多难猜''，不适合度量``有多大''——后者是方差的活。

第三，也是最容易被跳过的一条：微分熵不是离散熵的推广。它可以为负，会随坐标尺度改变，把长度单位从米换成毫米，数字就跳了 \(\log_2 1000\)。因此单个连续变量的``信息量''没有绝对意义，两列连续特征的微分熵不能直接比大小。有意义的是差——互信息与 KL 散度在坐标变换下不变，特征筛选该用它们。

第四，最大熵是一条建模纪律：在给定约束下选最不承诺的分布。给定支撑得到均匀分布，加上均值得到指数分布，再加上方差得到 Gaussian。这解释了这些分布为何总在``默认''位置出现，也划出了默认失效的场合——重尾数据上，二阶矩约束可能既估不准也不存在。

\subsection{本章篇目}

四篇按依赖顺序排列，从具体困惑进入也可以直接翻到最近的一篇，再沿篇内链接回补。

\begin{itemize}
\tightlist
\item \hyperref[sec:07-Information-Theory/01-Information-and-Entropy/01]{``从惊讶到 Bit 与自信息''}：为什么信息量必须是概率的负对数，底数怎样决定单位，以及提问树上那棵图；
\item \hyperref[sec:07-Information-Theory/01-Information-and-Entropy/02]{``Shannon 熵与二元熵''}：平均自信息、二元熵曲线、困惑度换算与压缩下界；
\item \hyperref[sec:07-Information-Theory/01-Information-and-Entropy/03]{``微分熵为何不是离散熵的直接复制''}：负值、尺度依赖、量化桥，以及哪些量仍然可信；
\item \hyperref[sec:07-Information-Theory/01-Information-and-Entropy/04]{``最大熵原理：从约束到分布''}：从约束反推分布，指数族与 softmax，以及重尾场合的失效。
\end{itemize}

只想拿走随手可用的直觉的话，记住三句就够：自信息衡量一次结果有多意外，熵是这种意外在模型下的平均，最大熵是在明确约束下不擅自补充结构。若要继续读编码、互信息与 KL，第三篇的坐标依赖值得多花点时间——后面所有连续情形的信息量，都建立在``只有差有意义''这一条上。

最后提醒一件事，它会在整章反复被违反：熵、微分熵、最大熵解，全都是相对于某个明确的概率模型而言的。变量是谁、分布由谁给出、对数取什么底、连续情形的参考测度是什么——这四项交代不清，算出来的数就没有完整含义。
